% No preamble. Prefix acc-. Fixed future classical system, no drift sign.
\section{A coupling coordinate with exact classical transport}
\label{acc-section}

Fix the countable joint infinite homogeneous classical system and
the invariant quotient recurrence constructed in Notebook30.
This section changes the marginal coordinate on that system.
It does not change the physical fields, the Gaussian windows or
the coupling used to define each actual local integral.
The construction uses the entire fixed future classical system;
no independence of different regulator subsequences is asserted.

\subsection{Uniform classical blocks and their marginal leakage}
\label{acc-input-section}

At stage $j$, let $\mathcal Q_j$ be the invariant analytic
interaction quotient with the same radius $b$ and carrier
weight in its current units. Let $S_j$ be the actual full
ordinary reference and $\beta_j$ the original Maxwell
functional. Set
\begin{equation}
 \beta_j(S_j)=1,\qquad
 P_j=I-S_j\beta_j,\qquad
 X_j=\ker\beta_j,\qquad
 A_j=P_{j+1}\mathcal T_j\big|_{X_j}.
 \label{acc-definitions}
\end{equation}
Here $\mathcal T_j$ is actual classical background
substitution, and the ordinary minimum identity gives
\begin{equation}
                         \mathcal T_jS_j=S_{j+1}.
 \label{acc-reference-identity}
\end{equation}
Recall the uniform bounds
\begin{equation}
 \|\beta_j\|\le C_\beta b^{-2},\qquad
 \|S_j\|\le C_Sb^2,\qquad
 \|P_j\|\le p,\qquad
 \|\mathcal T_j\|\le B_T.
 \label{acc-norm-inputs}
\end{equation}
The constants are independent of the later inner-radius
restriction and the homogeneous stage. They may depend
on the fixed geometry and support weights.

Assume, as supplied by the finite-block classical estimate,
that for one fixed integer $q\ge1$ and $0\le\theta<1$,
\begin{equation}
 \|A_j\|\le B,\quad B\ge1,\qquad
 \|A_{j+q-1}\cdots A_j\|_{X_j\to X_{j+q}}\le\theta
                         \quad\hbox{for every }j.
 \label{acc-block-input}
\end{equation}
Define
\begin{equation}
 U_{j,0}=I_{X_j},\qquad
 U_{j,n}=A_{j+n-1}\cdots A_j,\qquad
 \ell_j(e)=\beta_{j+1}(\mathcal T_je),\quad e\in X_j.
 \label{acc-evolution-leakage}
\end{equation}
No identity $\ell_j=0$ is assumed. Its norm satisfies
\begin{equation}
                         \|\ell_j\|\le C_\beta B_Tb^{-2}.
 \label{acc-leakage-bound}
\end{equation}
Writing $n=mq+r$ with $0\le r<q$ and grouping consecutive
blocks from the initial stage gives
\begin{equation}
            \|U_{j,n}\|\le B^r\theta^m
                         \le B^{q-1}\theta^{\lfloor n/q\rfloor}.
 \label{acc-evolution-decay}
\end{equation}
No commutation of different stages is used.

\subsection{A convergent dual correction}
\label{acc-dual-section}

\begin{theorem}[Adapted marginal functional]
\label{acc-adapted-theorem}
There is a unique uniformly bounded family of linear
functionals $k_j:X_j\to\mathbb C$ satisfying
\begin{equation}
                         k_j=\ell_j+k_{j+1}A_j.
 \label{acc-cohomological}
\end{equation}
It is given by the absolutely convergent dual-norm series
\begin{equation}
                k_j(e)=\sum_{n\ge0}\ell_{j+n}(U_{j,n}e),
 \qquad
 \|k_j\|\le C_kb^{-2},\qquad
 C_k=\frac{C_\beta B_T B^{q-1}q}{1-\theta}.
 \label{acc-series}
\end{equation}
The new functional
\begin{equation}
                  \widehat\beta_j=\beta_j+k_jP_j
 \label{acc-adapted-beta}
\end{equation}
satisfies
\begin{equation}
 \widehat\beta_j(S_j)=1,\qquad
 \|\widehat\beta_j\|\le C_{\widehat\beta}b^{-2},
 \quad C_{\widehat\beta}=C_\beta+C_kp,\qquad
 \widehat\beta_{j+1}\mathcal T_j=\widehat\beta_j.
 \label{acc-exact-transport}
\end{equation}
Uniqueness concerns the specified fixed future classical
system, not a comparison of different classical limits.
\end{theorem}
\begin{proof}
Equations \eqref{acc-leakage-bound} and
\eqref{acc-evolution-decay} give
\[
 \sum_{n\ge0}\|\ell_{j+n}U_{j,n}\|
 \le C_\beta B_Tb^{-2}B^{q-1}
                 \sum_{m\ge0}\sum_{r=0}^{q-1}\theta^m
 =C_kb^{-2}.
\]
Thus \eqref{acc-series} converges absolutely in the dual
Banach space. Separating its first term and shifting the
remaining index proves \eqref{acc-cohomological}.

If two uniformly bounded families solve that equation,
their difference $d_j$ obeys
$d_j=d_{j+n}U_{j,n}$ for every $n$.
The common bound and \eqref{acc-evolution-decay} force
$d_j=0$. This argument is for the fixed $b$ and the
fixed future system.

The normalization and norm in \eqref{acc-exact-transport}
follow immediately from $P_jS_j=0$ and
\eqref{acc-norm-inputs}. To prove the last identity,
decompose $x=sS_j+e$ with $s=\beta_j(x)$ and $e=P_jx$.
Then
\begin{equation}
 \mathcal T_jx=(s+\ell_j(e))S_{j+1}+A_je.
 \label{acc-classical-decomposition}
\end{equation}
Applying $\widehat\beta_{j+1}$ and using
\eqref{acc-cohomological} gives
$s+\ell_j(e)+k_{j+1}(A_je)=s+k_j(e)
=\widehat\beta_j(x)$.
\end{proof}

If the original data preserve the real form, every term
of \eqref{acc-series} is real on real inputs; the adapted
functionals therefore preserve that real form.
The new full-reference projection
\begin{equation}
 \widehat P_j=I-S_j\widehat\beta_j,\qquad
                 \|\widehat P_j\|\le1+C_SC_{\widehat\beta}
 \label{acc-adapted-projection}
\end{equation}
has a bound independent of $b$. The original $\beta_j$
remains the local Maxwell coefficient. The adapted
$\widehat\beta_j$ is a bounded coordinate choice, not a
newly asserted Ward classification.

For completeness, the two irrelevant spaces are explicitly
conjugate. Define
\begin{equation}
 J_j:X_j\longrightarrow\ker\widehat\beta_j,\qquad
                       J_j(e)=e-k_j(e)S_j.
 \label{acc-conjugating-map}
\end{equation}
Then
\begin{equation}
 \|J_j\|\le1+C_SC_k,\qquad
 J_j^{-1}=P_j\big|_{\ker\widehat\beta_j},\qquad
                    \mathcal T_jJ_j=J_{j+1}A_j.
 \label{acc-conjugacy}
\end{equation}
Indeed $P_jJ_j(e)=e$ and
$\widehat\beta_j(J_j(e))=0$. Conversely,
$\widehat\beta_j(x)=0$ implies
$x=P_jx-k_j(P_jx)S_j$. Finally
\eqref{acc-classical-decomposition} and
\eqref{acc-cohomological} give the intertwining identity
in \eqref{acc-conjugacy}. This does not claim a new
one-step contraction in an unchanged norm.

\subsection{The exact density coordinate}
\label{acc-density-section}

Write the normalized density at stage $j$ as
\[
                  \exp[-\gamma_jS_j+e_j],
            \qquad \gamma_j=g_j^{-2},\quad e_j\in X_j.
\]
Set
\begin{equation}
 \widehat\gamma_j=\gamma_j-k_j(e_j),\qquad
                  \widehat e_j=e_j-k_j(e_j)S_j.
 \label{acc-density-coordinate}
\end{equation}
Then $\widehat e_j\in\ker\widehat\beta_j$ and
\begin{equation}
                 -\widehat\gamma_jS_j+\widehat e_j
                              =-\gamma_jS_j+e_j
 \label{acc-density-equality}
\end{equation}
exactly. On the established quadratic envelope
$\|e_j\|\le C_{\rm int}b^2$,
\begin{equation}
 |\widehat\gamma_j-\gamma_j|\le C_kC_{\rm int},\qquad
 \|\widehat e_j\|\le(1+C_SC_k)C_{\rm int}b^2.
 \label{acc-coordinate-size}
\end{equation}
For real data a sufficiently large positive inverse coupling
therefore remains positive under this coordinate change.
The actual Gaussian map still uses the original
$g_j=\gamma_j^{-1/2}$ and its prescribed threshold.
Reevaluating it at $\widehat\gamma_j$ would change the
integration data and is not part of this identity.

Let the actual local quantum correction at that fixed
coupling be
\begin{equation}
             \mathcal F_{j,g}(e)
                   =\mathcal T_je+\mathcal G_{j,g}(e).
 \label{acc-quantum-decomposition}
\end{equation}
The original density-preserving extraction is
\begin{equation}
 \gamma_{j+1}
   =\gamma_j-\beta_{j+1}(\mathcal F_{j,g_j}(e_j)),\qquad
 e_{j+1}=P_{j+1}\mathcal F_{j,g_j}(e_j).
 \label{acc-original-update}
\end{equation}
Combining \eqref{acc-adapted-beta},
\eqref{acc-exact-transport} and
\eqref{acc-density-coordinate} yields the exact adapted drift:
\begin{align}
 \widehat\gamma_{j+1}
 &=\gamma_j-\widehat\beta_{j+1}
                           (\mathcal F_{j,g_j}(e_j))\notag\\
 &=\widehat\gamma_j-
                         \widehat\beta_{j+1}
                           (\mathcal G_{j,g_j}(e_j)).
 \label{acc-adapted-update}
\end{align}
The classical leakage has disappeared from the scalar
equation by a bounded change of coordinate. No sign
information has been introduced.

\subsection{The inherited contribution is small}
\label{acc-inherited-section}

Restate the fixed-coupling quantum quotient estimate from
Notebook30:
\begin{equation}
 \|\mathcal G_{j,g}(x)-\mathcal G_{j,g}(y)\|_{\mathcal Q,b}
 \le C_{\rm quant}\frac{gT(g)}b\,
                         \|x-y\|_{\mathcal Q,b}
 \label{acc-quantum-input-bound}
\end{equation}
on its common admissible neighborhood. It does not compare
two different coupling parameters. On the quadratic
amplitude envelope, \eqref{acc-exact-transport} gives
\begin{equation}
 \left|\widehat\beta_{j+1}
       (\mathcal G_{j,g}(e_j)-\mathcal G_{j,g}(0))\right|
 \le C_{\widehat\beta}C_{\rm quant}C_{\rm int}
                                      \frac{gT(g)}b.
 \label{acc-inherited-drift}
\end{equation}
For fixed positive $b$ and the prescribed thresholds,
the right side tends to zero as $g\downarrow0$.
Thus the leading adapted drift is insensitive to the
inherited irrelevant amplitude up to a controlled
vanishing error. The original quantity
$\beta_{j+1}(\mathcal T_je_j)$ need not itself be small.

\subsection{The remaining determinant and pivot coefficient}
\label{acc-leading-section}

For the step from $j$ to $j+1$ let
\begin{equation}
                 D_j=\Delta_j+\mathcal Q_{{\rm det},j}
                    \quad\hbox{as a class in }\mathcal Q_{j+1}.
 \label{acc-leading-definition}
\end{equation}
The fixed-threshold parity theorem of
Section~\ref{cgp-section}, evaluated afterward at $T(g)$,
gives on one fixed outer field radius $\rho$
\begin{equation}
 \|[\mathcal G_{j,g}(0)-D_j]\|_{\mathcal Q,\rho}
               \le C_{\rm par}\,[\tau(T(g))+g^2T(g)^6].
 \label{acc-outer-error}
\end{equation}
The constants are uniform in $j$ on the fixed homogeneous
system. Since the complete zero-input totals are invariant,
the norm limit $D_j$ is invariant as a total. For the
semisimple color group, the total first derivative of
the difference vanishes.

It is important to restrict this error correctly before
applying the inner-radius functional $\widehat\beta$.
For $b\le\rho/2$, remove the degree-zero and degree-one
background sectors coefficientwise and use homogeneous
phase projection. The normalized class is unchanged, and
the resulting bound is
\begin{equation}
 \|[\mathcal G_{j,g}(0)-D_j]\|_{\mathcal Q,b}
 \le \frac{2C_{\rm par}}{\rho^2}b^2
                     [\tau(T(g))+g^2T(g)^6].
 \label{acc-inner-error}
\end{equation}
This is precisely the inner-radius consequence proved in
\eqref{cgp-inner-bound}. Coefficientwise gauge invariance
is not required. The $b^2$ cancels the $b^{-2}$ in
$\|\widehat\beta\|$.

Define the remaining leading scalar coefficient by
\begin{equation}
                         b_{0,j}
                     =\widehat\beta_{j+1}(D_j).
 \label{acc-leading-coefficient}
\end{equation}
Equations \eqref{acc-adapted-update},
\eqref{acc-inherited-drift} and \eqref{acc-inner-error}
give the quantitative formula
\begin{align}
 &\left|
 \widehat\gamma_j-\widehat\gamma_{j+1}-b_{0,j}
 \right| \notag\\
 &\quad\le
 \frac{2C_{\widehat\beta}C_{\rm par}}{\rho^2}
                    [\tau(T(g_j))+g_j^2T(g_j)^6]
 +C_{\widehat\beta}C_{\rm quant}C_{\rm int}
                                      \frac{g_jT(g_j)}b.
 \label{acc-leading-drift}
\end{align}
All quantities use the actual physical $g_j$.
The inner radius is fixed before the weak-coupling limit.
The coefficient $b_{0,j}$ may depend on the chosen
classical system and on the stage. No stage limit,
numerical value or sign is asserted.

\subsection{What remains after this coordinate change}
\label{acc-scope-section}

There is a simpler special case. If an additional theorem for
the actual classical system proves
$\beta_{j+1}\mathcal T_j=\beta_j$, then $\ell_j=0$,
the series \eqref{acc-series} gives $k_j=0$, and
$\widehat\beta_j=\beta_j$ and
$\widehat\gamma_j=\gamma_j$. The present general argument
does not assume that identity and does not rule it out.
In that case \eqref{acc-leading-coefficient} reduces to
the ordinary Maxwell coefficient $\beta_{j+1}(D_j)$.

The future series \eqref{acc-series} defines the unique
bounded correction for one specified infinite classical
system. It is not a new locality assertion in RG time
and does not identify coordinates from unrelated
background subsequences. The construction is a bounded
change of the decomposition of an exact density, not a
physical field transformation.

The result reduces the leading scalar drift to the
concrete determinant/pivot coefficient in
\eqref{acc-leading-coefficient}, with the controlled
error \eqref{acc-leading-drift}. A suitable positive
lower bound on that coefficient, or a proved positive
block drift with the actual normalizations, would be
additional information. None is supplied by a norm
estimate or by the coordinate change.

In particular there is no contraction in the joint
coupling--activity variables, no proof that a physical
running trajectory stays weak indefinitely, and no
identification of the local sector with the full Wilson
measure. Uniform actual finite-regulator histories,
vacuum factors, large and mixed field sectors, and
continuum observables retain their separate obligations.
